\section{Introduction}

Reinforcement learning with verifiable rewards (RLVR) has become the default
recipe for improving symbolic reasoning: sample programs, execute them, reward
the correct ones, optimize with a policy-gradient method such as
GRPO~\cite{grpo2024}. Reported gains are large and consistent~\cite{deepseekr1}.

What those gains are measured against is remarkably uniform. Across the
empirical RLVR literature the control group consists of \emph{learning} arms:
the un-tuned base model, a matched supervised fine-tuning (SFT) run on the same
demonstrations, and---in the more careful papers---rejection-sampling
SFT~\cite{rft2023,restem2023}. Every one of these consumes the demonstration
budget and improves as it grows. A fourth option is almost never present: a
deterministic, non-learning symbolic solver. For program-annotated numerical
reasoning benchmarks such as FinQA~\cite{finqa2021} or
TAT-QA~\cite{tatqa2021}, where answers are short arithmetic programs over a
table, such a solver is not exotic---it is the classical approach the neural
pipeline replaced.

This paper asks whether that omission changes conclusions, and finds something
more general than expected.

\paragraph{The observation that starts it.}
With two demonstrations on CodeTAT-QA, warm-start SFT reaches $19.0\%$ dev
accuracy and GRPO reaches $72.0\%$---a 53-point gain that, against a
warm-start-only control, is an unambiguous success. A rule-based solver handed
the \emph{same} two demonstrations---and using neither---also reaches $72.0\%$.
On the frozen test
set the three seeds give $\Delta = -2.13$, $-1.42$, $-3.90$ points. The gain
over the learning baseline is real; the conclusion drawn from it is not.

\paragraph{Design.}
We pre-registered---before running anything on the second task---a single
decision quantity:
\begin{equation}
\label{eq:delta}
\Delta(m) = \mathrm{Acc}(\mathrm{GRPO}|W_m) - \max_{a \in \mathcal{B}} \mathrm{Acc}(a),
\end{equation}
with $\mathcal{B} = \{W_m, \text{c-SFT}, \text{RS-SFT}, \text{solver}\}$, all
arms branching from the same warm start $W_m$ trained on the same $m$
demonstrations. Taking a $\max$ over the control group, rather than comparing
against one chosen opponent, is the central design decision: it makes the
comparison impossible to ease by dropping an arm.

\paragraph{Findings.}
We evaluate two tasks and three backbones spanning $14\times$ in scale. All 17
frozen-test groups give negative $\Delta$, from $-1.06$ to $-16.42$ points
(Figure~\ref{fig:delta}). But the interesting structure is in \emph{which} arm
wins:

\begin{itemize}\itemsep2pt
\item At \textbf{0.5B and 1B}, the maximum is attained by the deterministic
solver in every one of 11 groups. Dose-response curves explain why: the solver
is flat in the demonstration budget while every learning arm scales steeply
with it, so the scarce regime is where a non-learning baseline is strongest.
\item At \textbf{7B}, the solver drops 15 points behind and is no longer
competitive---exactly as the ``RLVR gains emerge with backbone capability''
account predicts. $\Delta$ stays negative in all six groups anyway, because the
winning arm becomes \textbf{continued SFT}: the most elementary baseline in the
set, and one that appears in essentially every RLVR paper.
\end{itemize}

So the finding is not ``a symbolic solver beats RL''---that is a fact about one
scale regime. It is that \emph{at every scale we measured, some simpler method
matched or beat GRPO, but which method that is changes with scale}. A control
group holding only the base model and one matched SFT run misses the solver at
0.5B; a control group designed around the solver misses continued SFT at 7B.
What the evidence demands is breadth of the control group and a $\max$
aggregation over it, not the presence of any particular arm.

\paragraph{Contributions.}
(i) A control-group design for RLVR evaluation---a non-learning arm plus $\max$
aggregation---together with the pre-registered decision rule built on it;
(ii) an execution over two tasks and three backbones in which all 17
frozen-test effect sizes are negative and the identity of the winning baseline
shifts with scale; (iii) three independent observations of dev-selection bias,
each a non-negative $\Delta$ on dev that turns negative on the frozen test.

Section~\ref{sec:statistics} additionally audits the decision rule against the
selection bias it could itself introduce, and finds that it does not account for
the result. We report that as an audit of our own instrument rather than as a
contribution, alongside the per-group statistics---two of seventeen groups
significant under Holm---that the effect sizes actually support.

\paragraph{What we do not claim.}
We do not claim RLVR is ineffective: our GRPO runs beat every learning baseline
at 0.5B, by 35 points at $m{=}2$. We do not claim the solver generalizes---it is
task-specific and its engineering cost is reported in
Section~\ref{subsec:solvercost}. We do not claim the 7B groups establish a
scale effect: at ten times the prompt coverage the smallest-margin 7B cell
reverses to $\Delta = +6.74$ (Appendix~\ref{app:followup}), so Config~3 marks a
budget boundary rather than a capability one. We claim only that an RLVR
conclusion established without a broad control group, a stated budget and a
frozen test set is not trustworthy.
